\documentclass[12pt,a4paper]{article}
\usepackage{fontspec}
\defaultfontfeatures{Ligatures=TeX,Scale=MatchLowercase}
\IfFontExistsTF{Alice-Regular}{
  \setmainfont{Alice-Regular}[
      BoldFont = Alice-Regular,
      ItalicFont = Alice-Regular,
      BoldItalicFont = Alice-Regular,
      BoldFeatures = {FakeBold=1.6},
      ItalicFeatures = {FakeSlant=0.2},
      BoldItalicFeatures = {FakeBold=1.6,FakeSlant=0.2}
  ]
}{
  \setmainfont{Noto Serif}
}
\IfFontExistsTF{Noto Sans}{\setsansfont{Noto Sans}}{}
\usepackage{polyglossia}
\setmainlanguage{russian}
\usepackage[a4paper,left=20mm,right=20mm,top=20mm,bottom=20mm]{geometry}
\usepackage{graphicx}
\usepackage{amsmath,amssymb,amsthm}
\usepackage{booktabs,longtable,tabularx,array,float,multirow}
\usepackage{tikz}
\usepackage[europeanresistors]{circuitikz}
\usepackage{pgfplots}
\usepackage{xcolor}
\usepackage{caption}
\usepackage{fancyhdr}
\usepackage{enumitem}
\usepackage{tabularx}
\pgfplotsset{compat=1.18}
\pagestyle{fancy}
\fancyhf{}
\renewcommand{\headrulewidth}{0.4pt}
\renewcommand{\footrulewidth}{0.4pt}
\fancyfoot[C]{\thepage}
\newcolumntype{C}{>{\centering\arraybackslash}X}
\renewcommand{\arraystretch}{1.25}

\newcommand{\TitlePage}[2]{
\title{
	МИНОБРНАУКИ РОССИИ

	федеральное государственное автономное образовательное учреждение

	высшего образования

	«Национальный исследовательский университет

	«Московский институт электронной техники»
	\vspace{2em}
}
\author{}
\date{}
\maketitle
\thispagestyle{fancy}
\begin{center}
\Large Отчёт по лабораторной работе №#1\\
по дисциплине «Электротехника»\\
«#2»
\end{center}
\begin{center}
\Large Вариант 2
\end{center}
\vspace{12em}
\begin{table}[H]
\noindent
\begin{tabular*}{\textwidth}{@{\extracolsep{\fill}} p{0.48\textwidth} p{0.52\textwidth}}
Студент & Анисимов А.\,С., ЭН-$25$ \\
[0.5em] Руководитель & Самохин Д.\,В.
\end{tabular*}
\end{table}
\cfoot{Москва 2026}
\pagebreak
\fancyhead[L]{#2}
\fancyhead[R]{Электротехника}
}

\begin{document}
\TitlePage{2}{Исследование электрической цепи синусоидального тока}

\section*{Цель работы}
Исследовать основные характеристики синусоидального сигнала и режимы работы линейной цепи при чисто активной, индуктивной и ёмкостной нагрузке.

\section*{Исходные данные}
\begin{itemize}[leftmargin=1.2cm]
\item Для задания 1: $U_m=1.41\,\text{В}$, $f=60\,\text{Гц}$, $\psi_e=60^\circ$.
\item Для заданий 2--5: $u(t)=2.82\sin(3140t-30^\circ)\,\text{В}$.
\item В заданиях 3--5 использована схема методички с идеализированным измерением; сопротивлением шунта $0.01\,\Omega$ по сравнению с основным элементом пренебрегаем.
\end{itemize}

\section*{1. Характеристики синусоидального источника}
Действующее значение, период и угловая частота:
\[
U_{\text{RMS}}=\frac{U_m}{\sqrt2}=\frac{1.41}{\sqrt2}\approx 0.997\,\text{В},
\qquad
T=\frac1f=\frac1{60}\approx 0.01667\,\text{с},
\qquad
\omega=2\pi f=2\pi\cdot 60\approx 376.991\,\text{рад/с}.
\]

Следовательно,
\[
u(t)=1.41\sin(376.99t+60^\circ)\,\text{В}.
\]
Комплекс действующего значения:
\[
  U = U_{\text{RMS}}e^{j\psi_e} \approx 1.00e^{j60^\circ}\,\text{В}.
\]

\section*{2. Определение параметров сигнала по аналитическому выражению}
Для
\[
u(t)=2.82\sin(3140t-30^\circ)\,\text{В}
\]
получаем:
\[
U_m=2.82\,\text{В}, \qquad
U_{\text{RMS}}=\frac{2.82}{\sqrt2}\approx 1.994\,\text{В},
\]
\[
V_{p-p}=2U_m=5.64\,\text{В}, \qquad
f=\frac{\omega}{2\pi}=\frac{3140}{2\pi}\approx 499.747\,\text{Гц},
\qquad
T=\frac1f\approx 0.002001\,\text{с}.
\]
Начальная фаза сигнала:
\[
\psi_e=-30^\circ.
\]
Комплекс действующего значения:
\[
  U \approx 2.00e^{-j30^\circ}\,\text{В}.
\]

\section*{3. Активный элемент $R$}
Примем $R=1\,\text{к}\Omega$. Тогда
\[
Z_R=R=1000\,\Omega, \qquad
  I_R=\frac{  U}{R}=\frac{2e^{-j30^\circ}}{1000}
=2.00\cdot 10^{-3}e^{-j30^\circ}\,\text{А}.
\]

Мгновенное значение тока:
\[
i_R(t)=\frac{u(t)}R=\frac{2.82}{1000}\sin(3140t-30^\circ)
=2.82\cdot 10^{-3}\sin(3140t-30^\circ)\,\text{А}.
\]

Так как элемент чисто активный, ток и напряжение совпадают по фазе:
\[
\psi_i=\psi_u=-30^\circ.
\]

Средняя мощность на сопротивлении:
\[
P_R=U_{\text{RMS}}I_{\text{RMS}}
=1.994\cdot 0.001994
\approx 3.976\,\text{мВт}.
\]

Мгновенная мощность:
\[
p_R(t)=u(t)i(t)=U_mI_m\sin^2(3140t-30^\circ)
\approx 7.952\cdot 10^{-3}\sin^2(3140t-30^\circ)\,\text{Вт}.
\]

\section*{4. Индуктивный элемент $L$}
Примем $L=1\,\text{мГн}$. Индуктивное сопротивление:
\[
X_L=\omega L=3140\cdot 10^{-3}= 3.14\,\Omega.
\]
Тогда
\[
  I_L=\frac{  U}{jX_L}
=\frac{2e^{-j30^\circ}}{3.14e^{j90^\circ}}
\approx 0.635e^{-j120^\circ}\,\text{А}.
\]

Амплитуда тока:
\[
I_m=\frac{U_m}{X_L}=\frac{2.82}{3.14}\approx 0.898\,\text{А}.
\]
Следовательно,
\[
i_L(t)\approx 0.898\sin(3140t-120^\circ)\,\text{А}.
\]

Здесь ток отстаёт от напряжения на $90^\circ$, а средняя мощность равна нулю:
\[
  P_L = 0.
\]
Мгновенная мощность знакопеременна:
\[
p_L(t)= -\frac{U_mI_m}2\sin(6280t-60^\circ)
\approx -1.266\sin(6280t-60^\circ)\,\text{Вт}.
\]

Если подключить параллельно вторую такую же катушку, то
\[
L_{\text{экв}}=\frac{L\cdot L}{L+L}=\frac L2,
\]
поэтому $X_L$ уменьшается вдвое, а ток возрастает вдвое.

\section*{5. Ёмкостный элемент $C$}
Примем $C=1\,\text{нФ}$. Ёмкостное сопротивление:
\[
X_C=\frac1{\omega C}=\frac1{3140\cdot 10^{-9}}
\approx 318.471\,\text{к}\Omega.
\]
Тогда
\[
  I_C=\frac{  U}{-jX_C}
\approx 6.261\, \mu\text{А}\cdot e^{j60^\circ}.
\]

Амплитуда тока:
\[
I_m=\frac{U_m}{X_C}=\frac{2.82}{318471}
\approx 8.855\,\mu\text{А}.
\]
Следовательно,
\[
i_C(t)\approx 8.855\cdot 10^{-6}\sin(3140t+60^\circ)\,\text{А}.
\]

Для конденсатора ток опережает напряжение на $90^\circ$, а средняя мощность также равна нулю:
\[
  P_C=0.
\]
Мгновенная мощность:
\[
p_C(t)= \frac{U_mI_m}2\sin(6280t-60^\circ)
\approx 1.248527e-05\sin(6280t-60^\circ)\,\text{Вт}.
\]

Если параллельно подключить второй такой же конденсатор, то
\[
C_{\text{экв}}=C+C=2C,
\]
поэтому $X_C$ уменьшится вдвое, а ток возрастёт вдвое.

\section*{6. Сводная таблица результатов}
\begin{center}
\begin{tabularx}{\textwidth}{|p{2.0cm}|p{2.5cm}|C|p{2.8cm}|C|}
\hline
Элемент & \centering Импеданс & Комплекс тока & \centering Фазовый сдвиг & Средняя мощность \\
\hline
$R=1\,\text{к}\Omega$ &
\centering $1000\,\Omega$ &
$2.00\,\text{мА}\cdot e^{-j30^\circ}$ &
\centering $0^\circ$ &
$3.976\,\text{мВт}$ \\
\hline
$L=1\,\text{мГн}$ &
\centering $j\,3.14\,\Omega$ &
$0.635\,\text{А}\cdot e^{-j120^\circ}$ &
\centering $-90^\circ$ &
$0$ \\
\hline
$C=1\,\text{нФ}$ &
\centering $-j\,318.47\,\text{к}\Omega$ &
$6.261\,\mu\text{А}\cdot e^{j60^\circ}$ &
\centering $+90^\circ$ &
$0$ \\
\hline
\end{tabularx}
\end{center}

\section*{Вывод}
Рассчитаны все требуемые параметры синусоидального сигнала и режимы работы 
цепи с активной, индуктивной и ёмкостной нагрузкой. Показано, что в чисто активной цепи 
ток совпадает по фазе с напряжением и выделяется средняя активная мощность, а в идеальных 
реактивных элементах средняя мощность равна нулю, поскольку энергия только периодически 
накапливается и возвращается источнику. Дополнительно обосновано, почему при параллельном 
соединении одинаковых катушек или конденсаторов ток возрастает в два раза.

\end{document}
